\subsection{Validazione}
\label{subsec:proc_validazione}

Il processo di validazione ha lo scopo di accertare che il prodotto finale soddisfi i bisogni dell'utente e l'utilizzo di esso.
La validazione viene effettuata come controllo finale, e consiste nel verificare che il prodotto sia conforme alle attese e ai requisiti specificati nel documento \textit{Analisi dei Requisiti}.

\subsubsection{Implementazione della validazione}

Il raggiungimento della soddisfazione dei test (di unità, di integrazione e di sistema) descritti nel documento di \textit{Piano di Qualifica}, è un requisito necessario per poter procedere con la convalida del software.
La validazione deve essere fatta in presenza del referente dell'azienda proponente. Ciò include i seguenti passaggi:
\begin{itemize}
    \item \textbf{Revisione dei requisiti}: il gruppo controlla che tutti i requisiti siano stati soddisfatti;
    \item \textbf{Collaudo}: il gruppo esegue i test di accettazione per verificare che il prodotto implementi correttamente i requisiti e le funzionalità richieste.
\end{itemize}
Al termine di questi passaggi, il gruppo può procedere con la consegna del prodotto all'azienda proponente che lo convalida.

\subsubsection{Test di accettazione}
I test di accettazione sono dei test di tipo black-box, cioè si basano sull'esecuzione di specifici casi di test che vanno a riprodurre scenari d'uso realistici, replicando il comportamento degli utenti.
Questi vanno, inoltre, a verificare che il prodotto soddisfi i requisiti utente specificati nel documento di \textit{Analisi dei Requisiti}.
Poiché questi test mirano a simulare l'esperienza reale dell'utente, è fondamentale che vengano eseguiti su un prodotto completamente funzionante, in modo da garantire un'analisi accurata delle sue prestazioni in condizioni operative effettive.

\subsubsection{Strumenti}
Qui vengono elencati gli strumenti usati, che vengono poi approfonditi nella sezione \hyperref[sec:strumenti]{Strumenti}.
\begin{itemize}
    \item Meet;
    \item Zoom.
\end{itemize}